\chapter{問題設定}
% 【この章の目的】入力，出力，仮定，記号，最適化目標，評価対象を曖昧なく定義する．
% 内容入り例：examples/chapters/03-problem-example.tex

\section{対象と仮定}
% 対象データ，取得条件，利用可能な情報，適用範囲を記載する．

\section{記号と定式化}
% 主要記号は初出時に定義し，単位・次元・添字の意味を統一する．

\section{評価課題}
% 何をもって問題を解けたと判断するかを，実験章に先立って定義する．
